\scanpage{178}
\mspage{36}

\uline{Question} Si $X$ dans Cat est tot. $W_\infty$-asphérique, i.e. les $f_!(e_{X^\wedge})$ sont $W_\infty$-asphériques pour tte flèche $f:X\to Y$ de Cat, ou encore pour tte telle flèche, avec les $\,{}_{y}\backslash Y$. $0$-connexes ($y\in\operatorname{Ob}Y$), $Y$ est $W_\infty$-asphérique -- cela implique-t-il que $X$ soit \unsure{($'$)} filtrante ? Ça ne m'étonnerait pas, et $=$ \deleted{\unsure{en}} supposant seulement les $f_!(e_{X^\wedge})$ $1$-connexes$^*$ -- (ils sont $0$-connexes, dès que $X$ l'est). Mais je ne veux pas m'attarder à présent sur cette question\unsure{\ldots}

% slanted
\reversemarginpar\margin{\deleted{\uline{Non}, cf \unsure{plus bas}}}\normalmarginpar

% slanted
\reversemarginpar\margin{Ils le sont toujours \unsure{aussi|ainsi} (cf $\underline{\overline{\mathrm{XII}}}$ plutôt), cf \unsure{\S\,?} \unsure{c'est} \uline{\unsure{2-connexe}} qu'il faut lire \unsure{depuis tjrs|de suite} !}\normalmarginpar

La réponse à la question précédente est négative, comme on voit par la prop.

% slanted
\reversemarginpar\margin{\uuline{Faux} $\longrightarrow$ cf \unsure{déf. 3.3}, p.~45}\normalmarginpar

\uline{Prop.} 2.19. Soit $X$ \deleted{une catégorie} \inserted{de Cat}, telle que $\operatorname{diag}_X:X\to X\times X$ soit cofinale$^*$, i.e. telle que $X^{\unsure{o|0}}$ soit tot. $0$-connexe. Alors $X$ est tot. $\mathbb W_\infty$-asph. (donc aussi tot $W$-asphérique, pour tt \unsure{loc.} fond. \deleted{\unsure{tot}} tel que $W\supset\mathbb W_\infty$).

Soit en effet un foncteur cofinal $f:X\to Y$, il faut prouver que $Y$ est $\mathbb W_\infty$-asph. Or on sait que tot-$0$-connexe implique $\mathbb W_\infty$-asph.$^*$ (cf prop. 2.3 p.~17-18), donc il suffit de voir que $Y^o$ est tot. $0$-conn., i.e. $Y\xrightarrow{\deleted{\unsure{\to}}d_Y}Y\times Y$ cofinal. Or on a le diagramme commutatif

% slanted
\reversemarginpar\margin{C'est \uuline{faux}}\normalmarginpar

\diagram{\begin{tikzcd}[ampersand replacement=\&]
X\arrow[r,"d_X"]\arrow[d,"f"']\&X\times X\arrow[d,"f\times f"]\\
Y\arrow[r,"d_Y"']\&Y\times Y
\end{tikzcd}}

où par hypothèse $f$, donc aussi $f\times f$, et $d_X$ sont cofinaux. Donc $g=(f\times f)\circ\unsure{p_X|d_X}$ est cofinal, et comme $g=d_Yf$ et $f$ est aussi cofinal, on conclut que $d_Y$ l'est, qed.

\uline{Corollaire} $\mathord{\tikz[baseline=0pt,x=1em,y=1em]{\draw (0,0)--(.45,.72)--(.74,0)--cycle; \draw (.13,.025)--(.48,.63);}}^{\deleted{\unsure{o}}}$ \deleted{\unsure{étant}} tot. $0$-conn, $(\mathord{\tikz[baseline=0pt,x=1em,y=1em]{\draw (0,0)--(.45,.72)--(.74,0)--cycle; \draw (.13,.025)--(.48,.63);}})^o$ est \inserted{tot.} $\mathbb W_\infty$-asph. pourtant il n'est pas filtrant.

% slanted
\margin{Question ($X$ \deleted{\unsure{tot.}}$-\mathbb W_\infty$-asph $\Rightarrow X^o$ tot. $0$-connexe ? On dirait \unsure{oui} puisque \unsure{PS$'1$} \unsure{--}}

% archive sheet number: 178
